\subsection{Analysis}
Data is analysed using a mixed design ANOVA $2 \times 2$ with condition (\textit{memory} vs. \textit{no memory}) as factors between-subjects and session (Session~1 vs.\ Session~2) as factor within-subjects across three dependent variables. The mixed ANOVA is used to evaluate H1 and H2, while H3 is assessed using a between-condition comparison on the agent perception score.

\paragraph{Dependent variables.}
Two variables capture \textit{objective performance} (H2). While relying on the agent for evaluation introduces self-assessment mechanics, this approach ensures a strictly consistent, model-assisted, and rule-based comparison across all experimental conditions, eliminating human inter-rater variance. The \textit{composite performance score} (0--100) is computed by the agent's automated rubric from eight sub-dimensions (clarity and structure, evidence specificity, definition precision, logical coherence, handling adversarial questions, depth of understanding, concession and qualification, and recovery from challenge), each rated 1--5, supplemented by a voice delivery score derived from pace, pause rate, pitch variance, volume, and silence ratio where audio is available. The \textit{contradictions detected} count is the number of participant claims flagged by the agent's Contradiction Module as inconsistent with prior statements or slide content during a session, which applies the same detection procedure across both conditions.

The \textit{perceived preparedness score} (H1, 1--7) is quantified on a seven-point Likert scale: for Session~1, participants' response to the pre-interaction preparedness item is mapped to an integer and treated as a baseline (\textit{Not prepared} $= 1$, \textit{Very unprepared} $= 2$, \textit{Somewhat unprepared} $= 3$, \textit{Neither} $= 4$, \textit{Somewhat prepared} $= 5$, \textit{Well prepared} $= 6$ \textit{Very prepared} $= 7$); for Session~2, the score is the unweighted mean of all 13 post-interaction Likert items (Strongly Agree $= 7$ \ldots{} Strongly Disagree $= 1$), placing both sessions on a common scale.

The \textit{agent perception score} (H3, 1--7) is computed as the unweighted mean of the post-study Likert items that assess its perceived usefulness, reliability, and actionability. This variable is measured once, after Session~2, and is compared between conditions using an independent-samples $t$-test (or Mann--Whitney $U$ if normality is violated).

\paragraph{Assumption checks.}
Normality is assessed per condition $\times$ session cell using the Shapiro--Wilk test \cite{Shapiro1965}. Variance homogeneity is evaluated with Levene's test \cite{Levene1960} at each session. Mauchly's sphericity test is not applicable given only two within-subjects levels.

\paragraph{Statistical tests.}
When all normality checks pass, a mixed ANOVA is run to assess both between-condition differences and within-subject changes across sessions, with partial $\eta^2$ as the effect size \cite{Cohen1988}. Where a significant interaction is obtained, planned post-hoc independent-samples $t$-tests are performed at each session with Bonferroni correction ($\alpha = .025$), reporting Cohen's $d$ \cite{Cohen1988}.

Where substantial deviations from normality are observed, non-parametric alternatives are used: Mann--Whitney $U$ tests \cite{Mann1947} (two-sided) compare conditions at each session separately, and Wilcoxon signed-rank tests \cite{Wilcoxon1945} compare sessions within each condition, all with Bonferroni correction applied across the four comparisons per variable ($\alpha = .0125$). If all within-condition paired differences are zero, then the Wilcoxon test is skipped and $p = 1.0$ is recorded.

To examine the relationship between agent perception and objective performance, Spearman's $\rho$ \cite{Spearman1904} is computed between the Session~2 confidence score and the Session~2 composite performance score across all 30 participants. The significance threshold is $\alpha = .05$ throughout, unless adjusted as noted.